\section{Rangkuman}
\begin{itemize}
    \item ElGamal (Taher ElGamal, 1985) adalah sistem kunci-publik yang keamanannya bersandar pada persoalan logaritma diskrit, erat terkait dengan Diffie-Hellman.
    \item Pembangkitan kunci: pilih prima $p$ dan akar primitif $g$, pilih kunci privat $x$ dengan $2 \le x \le p-2$, lalu hitung kunci publik $y = g^x \bmod p$.
    \item Enkripsi bersifat probabilistik: pilih bilangan acak $k$, hitung $a = g^k \bmod p$ dan $b = y^k \cdot m \bmod p$, sehingga cipherteksnya berupa pasangan $(a, b)$.
    \item Dekripsi: $m = b \cdot (a^x)^{-1} \bmod p$, dengan $(a^x)^{-1} \equiv a^{p-1-x} \pmod p$.
    \item Contoh: $p = 2273$, $g = 3$, dan $x = 243$ memberi $y = 461$; pesan $m = 700$ dengan $k = 1463$ menjadi $(a, b) = (1439, 74)$ dan kembali menjadi $700$ setelah dekripsi.
    \item Karena $k$ berubah setiap sesi, plainteks yang sama menghasilkan cipherteks berbeda; hal ini mencegah serangan \textit{codebook} dan analisis statistik.
    \item Kelemahannya adalah ekspansi cipherteks (dua kali panjang plainteks); dibanding RSA, ElGamal bersifat probabilistik dan mendasari skema tanda tangan DSA.
\end{itemize}
